% PAGES: 136-138
% Continues sec04_c2_1.tex (pages 134-135). No environment was left open at
% the end of that file (the Summary bullet list and closing remark both
% close cleanly); this file opens fresh with Theorem 4.3, still under
% section 4.1 "Definitions and properties". Section 4.2 "Diagonal forms"
% begins partway down page 138 (folio 122).

\begin{theorem}
(i) The matrix $A$ is singular if and only if $0$ is an eigenvalue of $A$.

(ii) The matrix $A$ is nonsingular if and only if $0$ is not an eigenvalue of $A$.
\end{theorem}

\begin{proof}
(i) From (1.75), we find that the matrix $A$ is singular if and only if
\[
\Null(A) \;=\; \{v \in \R^n \ \text{such that} \ Av = 0\} \;\neq\; \{0\},
\]
i.e., if and only if there exists a vector $v_0 \neq 0$ such that $Av_0 = 0$, which is
equivalent to saying that $0$ is an eigenvalue of $A$; see (4.1).

(ii) Since a matrix is nonsingular if and only if the matrix is not singular, it
follows from (i) that a matrix is nonsingular if and only if $0$ is not an eigenvalue
of $A$.
\end{proof}

\begin{lemma}
Let $A$ be a nonsingular matrix, and let $A^{-1}$ be the inverse matrix of $A$. If
$(\lambda, v)$ are an eigenvalue and a corresponding eigenvector of $A$, then
$(\frac{1}{\lambda}, v)$ are an eigenvalue and a corresponding eigenvector of
$A^{-1}$.
\end{lemma}

\begin{proof}
If $v \neq 0$ is the eigenvector corresponding to the eigenvalue $\lambda$ of $A$,
then $Av = \lambda v$. Multiplying to the left by $A^{-1}$, we obtain that
$A^{-1}Av = \lambda(A^{-1}v)$, and, using the fact that $A^{-1}Av = v$ since
$A^{-1}A = I$, it follows that
\begin{equation}
v \;=\; \lambda\,(A^{-1}v).
\end{equation}
Note that $\lambda \neq 0$, since $A$ is a nonsingular matrix; cf.\ Theorem 4.3. Then,
from (4.16), we find that
\[
A^{-1}v \;=\; \frac{1}{\lambda}v.
\]
Since $v \neq 0$, it follows that $\frac{1}{\lambda}$ is an eigenvalue of $A^{-1}$
with corresponding eigenvector $v$.
\end{proof}

\begin{lemma}
Let $A$ be a square matrix, and let $\lambda$ and $v$ be an eigenvalue and a
corresponding eigenvector of $A$.

(i) Let $k$ be a positive integer. Then, $\lambda^k$ is an eigenvalue of the matrix
$A^k$ and $v$ is a corresponding eigenvector.

(ii) Let $P(x)$ be an arbitrary polynomial. Then, $P(\lambda)$ is an eigenvalue of the
matrix $P(A)$ and $v$ is a corresponding eigenvector.
\end{lemma}

\begin{proof}
(i) We show by induction over the positive integer $k \geq 1$ that
\begin{equation}
A^k v \;=\; \lambda^k v, \quad \forall\, k \geq 1.
\end{equation}
Formula (4.17) holds for $k=1$, since $\lambda$ and $v$ are an eigenvalue of $A$ and a
corresponding eigenvector, and therefore $Av = \lambda v$. Assume that (4.17) holds
for $k$, i.e., assume that $A^k v = \lambda^k v$. Then,
\[
A^{k+1}v \;=\; A(A^k v) \;=\; A(\lambda^k v) \;=\; \lambda^k \cdot Av \;=\;
\lambda^k \cdot \lambda v \;=\; \lambda^{k+1}v.
\]
Thus, formula (4.17) is established for $k+1$ and the proof by induction of (4.17) is
complete.

(ii) Let $P(x) = \sum_{k=0}^p c_k x^k$, with $c_p \neq 0$, where $p = \deg(P)$ is the
degree of the polynomial $P(x)$. Using (4.17), we obtain that
\begin{align*}
P(A)v &= \left( \sum_{k=0}^p c_k A^k \right) v \;=\; \sum_{k=0}^p c_k \left( A^k v
\right)\\
&= \sum_{k=0}^p c_k \left( \lambda^k v \right) \;=\; \left( \sum_{k=0}^p c_k
\lambda^k \right) \cdot v\\
&= P(\lambda) v.
\end{align*}
Thus, $P(A)v = P(\lambda)v$, and we conclude that $P(\lambda)$ is an eigenvalue of the
matrix $P(A)$ and $v$ is a corresponding eigenvector.
\end{proof}

\begin{lemma}
If $A$ is a square matrix, then the matrices $A$ and $A^t$ have the same eigenvalues.
\end{lemma}

\begin{proof}
Recall from Theorem 4.1 that the eigenvalues of a matrix are the roots of its
characteristic polynomial. Thus, the eigenvalues of the matrix $A$ are the roots of
$P_A(x) = \det(xI-A)$, and the eigenvalues of the matrix $A^t$ are the roots of
$P_{A^t}(x) = \det(xI-A^t)$.

Since $(xI-A)^t = xI-A^t$, and using the fact that $\det((xI-A)^t) = \det(xI-A)$, see
Lemma 10.2, we obtain that
\begin{align*}
P_A(x) &= \det(xI-A) \;=\; \det((xI-A)^t) \;=\; \det(xI-A^t)\\
&= P_{A^t}(x).
\end{align*}
Thus, the polynomials $P_A(t)$ and $P_{A^t}(t)$ have the same roots, and we conclude
that the matrices $A$ and $A^t$ have the same eigenvalues.
\end{proof}

\begin{definition}
Let $A$ be a square matrix of size $n$. The trace of the matrix $A$ is the sum of its
main diagonal entries, i.e.,
\begin{equation}
tr(A) \;=\; \sum_{i=1}^n A(i,i).
\end{equation}
\end{definition}

\begin{lemma}
Let $A$ be a square matrix of size $n$, let $P_A(t)$ be the characteristic polynomial
of $A$, and let $\lambda_i$, $i = 1:n$, be the eigenvalues of $A$. Then,
\begin{equation}
P_A(t) \;=\; t^n \;-\; tr(A)t^{n-1} \;+\; \dots \;+\; (-1)^n \det(A),
\end{equation}
and
\begin{equation}
\sum_{i=1}^n \lambda_i \;=\; tr(A);
\end{equation}
\begin{equation}
\prod_{i=1}^n \lambda_i \;=\; \det(A).
\end{equation}
\end{lemma}

\begin{proof}
Recall from Lemma 4.1 that
\begin{align}
\det(tI-A) &= \prod_{i=1}^n (t-\lambda_i)\\
&= t^n \;-\; \left( \sum_{i=1}^n \lambda_i \right) t^{n-1} \;+\; \dots \;+\;
(-1)^n \prod_{i=1}^n \lambda_i.
\end{align}
Let $t=0$ in (4.22--4.23). Then,
\[
\det(-A) \;=\; (-1)^n \det(A) \;=\; (-1)^n \prod_{i=1}^n \lambda_i,
\]
and therefore (4.21) is established.

Note that the only terms of order $n-1$ from $\det(tI-A)$ are obtained by multiplying
the main diagonal entries of $tI-A$. In other words, the coefficient of the term
$t^{n-1}$ in $\det(tI-A)$ is the same as the coefficient of the term $t^{n-1}$ in
$\prod_{i=1}^n (t-A(i,i))$, which is equal to $-\sum_{i=1}^n A(i,i) = -tr(A)$; see
Definition 4.5. Then, we conclude from (4.23) that $tr(A) = \sum_{i=1}^n \lambda_i$.
\end{proof}

\section{Diagonal forms}

\begin{definition}
Let $A$ be a square matrix of size $n$. The matrix $A$ has a diagonal form if and only
if there exists an $n \times n$ nonsingular matrix $V$ and an $n \times n$ diagonal
matrix $\Lambda$ such that
\begin{equation}
A \;=\; V\Lambda V^{-1}.
\end{equation}
A matrix that has a diagonal form is called a diagonalizable matrix.
\end{definition}

\begin{theorem}
Let $A$ be a diagonalizable matrix, and let $A = V\Lambda V^{-1}$ be the diagonal form
of $A$. Then, the diagonal entries of the matrix $\Lambda$ are the eigenvalues of $A$
and the columns of the matrix $V$ are the corresponding eigenvectors of $A$.

More precisely, if $\Lambda = diag(\lambda_k)_{k=1:n}$ and $V = col(v_k)_{k=1:n}$,
then $\lambda_k$, $k = 1:n$, are the eigenvalues of $A$, and $v_k$ is an eigenvector
corresponding to $\lambda_k$, for $k = 1:n$, respectively.
\end{theorem}

\begin{proof}
By multiplying $A = V\Lambda V^{-1}$ to the right by $V$, we obtain that
\begin{equation}
A = V\Lambda V^{-1} \iff AV = (V\Lambda V^{-1})V \iff AV = V\Lambda,
\end{equation}
since $V^{-1}V = I$. Recall from (1.11) and (1.88) that $AV = col(Av_k)_{k=1:n}$ and
$V\Lambda = col(\lambda_k v_k)_{k=1:n}$. Then, from (4.25), we find that
\begin{align*}
A = V\Lambda V^{-1} &\iff col(Av_k)_{k=1:n} = col(\lambda_k v_k)_{k=1:n}\\
&\iff Av_k = \lambda_k v_k, \quad \forall\, k = 1:n,
\end{align*}
i.e., $\lambda_k$ is an eigenvalue of $A$ and $v_k$ is a corresponding eigenvector,
for $k = 1:n$.

The vectors $v_k$, $k = 1:n$, are linearly independent vectors since they are the
columns of the nonsingular matrix $V$.
\end{proof}
